\documentclass[12pt]{article}
\usepackage[a4paper, margin=1in, landscape]{geometry}
\usepackage{amsmath}
\usepackage{booktabs}
\usepackage{xcolor}
\usepackage{colortbl}
\usepackage{array}
\usepackage{enumitem}

\definecolor{headerblue}{HTML}{1a2744}
\definecolor{altrow}{HTML}{eef4fc}

\begin{document}
\begin{center}
    {\Large \textbf{Computer Architecture -- Lab 4}}\\[0.5cm]
    Prithish Chaudhary\\[0.2cm]
    March 2026
\end{center}
\vspace{0.8cm}
\section*{Performance Statistics Summary}
\begin{table}[h]
    \centering
    \renewcommand{\arraystretch}{1.6}
    \begin{tabular}{>{\raggedright\arraybackslash}p{3cm}
                    >{\centering\arraybackslash}p{2.5cm}
                    >{\centering\arraybackslash}p{2cm}
                    >{\centering\arraybackslash}p{2cm}
                    >{\centering\arraybackslash}p{2cm}
                    >{\centering\arraybackslash}p{2.5cm}
                    >{\centering\arraybackslash}p{3.5cm}}
        \toprule
        \rowcolor{headerblue}
        \color{white}\textbf{Benchmark} & 
        \color{white}\textbf{Instructions} & 
        \color{white}\textbf{Cycles} & 
        \color{white}\textbf{IPC} & 
        \color{white}\textbf{CPI} & 
        \color{white}\textbf{OF Stalls (Data Hazard)} & 
        \color{white}\textbf{Wrong Branch Instructions}\\
        \midrule
        \rowcolor{altrow}
        descending.out & 189 & 443 & 0.427 & 2.344 & 74  & 176 \\
        evenorodd.out  & 6   & 16  & 0.375 & 2.667 & 6   & 0   \\
        \rowcolor{altrow}
        fibonacci.out  & 62  & 118 & 0.525 & 1.903 & 20  & 32  \\
        palindrome.out & 42  & 107 & 0.393 & 2.548 & 47  & 14  \\
        \rowcolor{altrow}
        prime.out      & 24  & 53  & 0.453 & 2.208 & 15  & 10  \\
        \bottomrule
    \end{tabular}
    \caption{Execution metrics for various benchmark programs.}
    \label{tab:metrics}
\end{table}

\section*{Observations}

All cycle counts satisfy the relation: $\text{Cycles} = \text{Instructions} + 4\;(\text{pipeline drain}) + \text{OF Stalls} + \text{Wrong Branch Instructions}$.

\begin{enumerate}[leftmargin=*]

\item \textbf{Descending (Bubble Sort)} -- This benchmark has the highest cycle count (443) and CPI (2.344). Being a nested-loop sorting algorithm, it executes the most instructions (189) and has the highest wrong-branch count (176). The frequent \texttt{blt}, \texttt{bgt}, and \texttt{jmp} instructions inside the inner loop cause many control hazard penalties, since every taken branch flushes 2 instructions from the pipeline. The 74 data hazard stalls come mainly from \texttt{load}$\rightarrow$\texttt{use} dependencies (loading array elements immediately before comparison).

\item \textbf{Even or Odd} -- The simplest benchmark with only 6 instructions and no loops. It has 0 wrong-branch instructions (the single \texttt{beq} is not taken for the input 11). The 6 OF stalls arise from two back-to-back data dependencies: \texttt{load}$\rightarrow$\texttt{divi} on \texttt{x3}, and \texttt{divi}$\rightarrow$\texttt{beq} on \texttt{x31}, and \texttt{sub}$\rightarrow$\texttt{addi} on \texttt{x10}. Each stalls for 2 cycles due to the 5-stage pipeline depth. Its CPI of 2.667 is the worst, because the stall overhead is large relative to the tiny instruction count.

\item \textbf{Fibonacci} -- This benchmark has the best IPC (0.525) and lowest CPI (1.903) among all benchmarks. The loop body (\texttt{add}, \texttt{store}, \texttt{subi}, \texttt{addi}, \texttt{add}, \texttt{add}, \texttt{jmp}) is 7 instructions long with relatively few data dependencies between consecutive instructions. The \texttt{store} instructions do not mark a destination register busy, which avoids stalls. The 32 wrong-branch instructions come from the unconditional \texttt{jmp} (always flushes 2 instructions) and the loop-exit \texttt{blt} being taken on each iteration.

\item \textbf{Palindrome} -- This benchmark has the second-highest CPI (2.548) and the highest stall-to-instruction ratio (47 stalls / 42 instructions $\approx$ 1.12). The main loop uses \texttt{divi} to extract digits, which creates a chain of dependencies: \texttt{divi} writes both its destination register and \texttt{x31} (remainder), and subsequent instructions immediately read both. The tight dependency chain \texttt{divi}$\rightarrow$\texttt{addi x31}$\rightarrow$\texttt{muli}$\rightarrow$\texttt{add}$\rightarrow$\texttt{divi} causes frequent stalls on nearly every iteration, making this the most stall-intensive benchmark per instruction.

\item \textbf{Prime} -- This benchmark checks primality via trial division. Its loop body contains \texttt{div} (which writes both the quotient register and \texttt{x31}), followed immediately by \texttt{beq x0, x31} to check the remainder. This creates a 2-cycle stall per iteration on the \texttt{x31} dependency. The 10 wrong-branch instructions come from the \texttt{jmp for} (unconditional, 2 flushes each) and the final taken \texttt{bgt} when exiting the loop.

\end{enumerate}

\noindent\textbf{General Trends:}
\begin{itemize}[leftmargin=*]
    \item Programs with tight loops and back-to-back register dependencies (palindrome, descending) suffer the most data hazard stalls, since the interlocking mechanism stalls the OF stage for up to 2 cycles per dependency.
    \item Wrong-branch penalties dominate in loop-heavy programs: each taken branch or unconditional jump flushes 2 incorrectly fetched instructions, accounting for the majority of the overhead in descending (176/443 $\approx$ 40\% of cycles) and fibonacci (32/118 $\approx$ 27\%).
    \item Benchmarks with longer loop bodies and fewer consecutive dependencies (fibonacci) achieve better IPC, since pipeline stalls are amortized over more useful instructions per iteration.
    \item The \texttt{div}/\texttt{divi} instructions are particularly expensive because they create dependencies on two registers (the explicit destination and the implicit \texttt{x31} for the remainder), leading to more stall opportunities than other ALU instructions.
\end{itemize}

\end{document}
